% ECONOMICS_PROBLEM: {"id": "D86-634", "title": "Long-term contracts versus spot markets", "jel_code": "D86", "source_id": "OP-0293"}
% !TeX program = xelatex
\documentclass[11pt,a4paper]{article}
\usepackage[margin=23mm]{geometry}
\usepackage{fontspec}
\setmainfont{Latin Modern Roman}
\usepackage{amsmath,amssymb,mathtools}
\usepackage{xcolor,enumitem,booktabs,longtable,array,xurl,hyperref}
\definecolor{muted}{HTML}{53636C}
\hypersetup{colorlinks=true,urlcolor=blue,linkcolor=blue}
\setlength{\parindent}{0pt}
\setlength{\parskip}{5pt}
\newcommand{\R}{\mathbb R}
\newcommand{\E}{\mathbb E}
\newcommand{\N}{\mathbb N}
\newcommand{\ind}{\mathbf 1}
\newcommand{\Var}{\operatorname{Var}}
\newcommand{\Cov}{\operatorname{Cov}}
\newcommand{\supp}{\operatorname{supp}}
\DeclareMathOperator*{\argmin}{arg\,min}
\DeclareMathOperator*{\argmax}{arg\,max}
\newcommand{\entrymeta}[3]{{\sffamily\small\color{muted}\textbf{\texttt{#1}}\quad|\quad #2\par #3\par}}
\newcommand{\Problemref}[1]{\texttt{#1}}
\newcommand{\JELref}[2]{\texttt{#2} (JEL \texttt{#1})}
\begin{document}
\section*{Long-term contracts versus spot markets}
\textbf{Problem ID:} \texttt{D86-634}\quad \textbf{JEL:} \texttt{D86}\par
% BEGIN ECONOMICS STATEMENT
\label{problem:OP-0293}
{\sffamily\small\color{muted}Original subject group: Mechanism design and market design\par}
\label{definitions:problem:30007575}
\textbf{Audit classification:} Published research agenda. The survey’s §8 explicitly asks why long-term arrangements are scarce. The causal estimands and experiment are editorial; universal superiority is not established.
\subsection*{Definitions and precise research target}
\textbf{Model and research target.} Consider repeated buyer–seller relationships over \(T\) periods, with observed demand states \(s_t\), private valuations, and measured enforcement and commitment costs. Randomly offer otherwise comparable groups either a prespecified long-term contract menu \(Z=1\) or a sequence of spot-market opportunities \(Z=0\). Let \(W_i(Z)\) be relationship \(i\)'s discounted realized total surplus net of implementation costs, including the buyer's value minus resource cost, and let \(D_i(Z)\) indicate adoption of the offered long-term contract. Elicit buyer and seller forecasts before assignment and define belief divergence \(H_i\) by a stated distance between their predictive distributions. Principal estimands are
\[
 \Delta(h,c)=E[W_i(1)-W_i(0)\mid H_i=h,C_i=c],
 \qquad E[D_i(1)\mid H_i=h,C_i=c],
\]
where \(C_i\) summarizes prespecified enforcement conditions. Identify contract-offer effects using randomization, distinguish adoption effects under explicit instrumental-variable assumptions, and estimate whether belief disagreement explains limited adoption despite potential surplus gains. Use relationship-level assignment, report noncompliance and exit, and include common uncertainty and measurement error in forecasts. Transporting results across platforms requires overlapping contract and state support.

\emph{Definitions and maintained assumptions (editorial unless a source definition is named).} A unit \(i\) is one buyer--seller relationship, with both parties included in its realized surplus. Define \(W_i(z)=\sum_{t=1}^T\beta^{t-1}[v_i(x_t(z),\theta_t(z))-c_t(x_t(z))-k_t(z)]\); transfers between the two parties cancel only within that two-party welfare objective. Assignment \(Z\) is randomized independently of baseline covariates and potential outcomes, with no cross-relationship interference or a stated cluster interference model. Adoption \(D_i(1)\) is whether the specified long-term contract is accepted; define \(D_i(0)\) if adoption effects are studied. \(H_i\) is a predetermined distance, e.g. total variation, between elicited probability forecasts on a common finite state partition. \(\Delta(h,c)\) is a regular conditional treatment effect, only identified on covariate support; exact equality conditioning on continuous \(H,C\) needs smoothness or bins for estimation. Random contract offers identify intention-to-treat surplus effects. Effects of adoption itself additionally require exclusion, relevance and monotonicity for an instrumental-variable interpretation. Report enforcement costs, default and exit in welfare.
\subsection*{Variants and assumption changes}
\begin{enumerate}
\item \textbf{Exogenous enforcement (editorial).} Randomize enforcement strength independently of contract availability and estimate how it modifies offer effects and adoption.
\item \textbf{Forecast disagreement (editorial).} Randomize common forecasting information before contract offers, holding contract terms fixed. Test whether disagreement predicts or mediates adoption while stating exclusions needed for a causal belief channel.
\end{enumerate}
\subsection*{References and source audit}
\begin{enumerate}
\item §8 printed p.49 in June2018 revision. Supports the stated research area; exact displayed specification is editorial. DOI publication is2019; inspected revision is June2018. \url{https://cowles.yale.edu/sites/default/files/2022-09/d2102-r.pdf}
\end{enumerate}
\begin{enumerate}\raggedright
\item\hypertarget{definitions:source:30007575:1}{} Dirk Bergemann; Juuso Välimäki. 2019. \emph{Dynamic Mechanism Design: An Introduction}. Published JEL57(2),235–274 (2019); inspected June19,2018 Cowles revision §8, printed pp.49–51.
\url{https://cowles.yale.edu/sites/default/files/2022-09/d2102-r.pdf}.
DOI: \url{https://doi.org/10.1257/jel.20180892}.
\end{enumerate}

\subsection*{Common conventions: Source mathematical conventions}

These conventions apply to each entry unless explicitly replaced.

\textbf{Models and identification.} A primitive model \(m\) specifies all probability laws, feasible choices, information, equilibrium and measurement rules used by its target. The admissible class \(\mathcal M\) is fixed independently of outcomes. If \(P_O(m)\) is its observable probability law and \(\tau(m)\) the target, its identified set at observable law \(P\) is
\[
 \mathcal I_\tau(P)=\{\tau(m):m\in\mathcal M,\ P_O(m)=P\}.
\]
Point identification means that this set is a singleton. An empty set signals incompatibility of data and assumptions. Coordinate bounds need not describe the joint set. Bounds are sharp only if every retained target is attainable or arbitrarily approachable by compatible models. Finite bounds require explicit restrictions on unobserved quantities; unrestricted confounding, hidden wealth, extrapolation or arbitrary equilibrium selection can yield uninformative sets.

\textbf{Causal interventions.} A potential outcome \(Y_i(p)\) is the outcome under a complete policy regime \(p\), including timing, financing, information and market spillovers. The target population and units are fixed before treatment. With interference, use the complete assignment vector or a justified exposure mapping; individual no-interference notation alone cannot identify an economy-wide intervention. A difference of mean potential outcomes does not require the cross-world joint distribution. Conditioning on post-treatment migration, survival, enrollment or employment changes the estimand and needs separate justification.

\textbf{Statistical statements.} An estimator, confidence set, forecast or test specifies a sampling law, model class and loss/coverage criterion. Uniform \((1-\alpha)\)-coverage means \(\inf_{m\in\mathcal M}\Pr_m\{\tau(m)\in C_n\}\geq1-\alpha\), or an explicitly asymptotic version. The whole parameter space trivially has coverage, so informativeness requires a stated width, risk or power objective. A forecasting improvement is relative to a named benchmark, information set and evaluation population.

\textbf{Equilibrium and optimization.} An equilibrium comprises feasible plans, optimality, beliefs where needed, budget constraints and all market/resource clearing. The primitives or a stated benchmark define each correspondence \(\mathcal E(p;m)\). Nonemptiness is assumed or proved, never inferred just from compact policy sets. A selected-equilibrium target uses a declared selection rule; a robust target quantifies over all permitted equilibria. Use suprema or infima unless attainment is established. An \(\varepsilon\)-optimal policy is feasible and within \(\varepsilon\) of the finite optimum value. Report an empty feasible set separately.

\textbf{Welfare and accounting.} Normative utility, interpersonal normalization, social weights, numeraire, discounting and the use of public receipts are primitives. Behavioral utility need not coincide with the welfare criterion. Household welfare includes ownership income and rents once; adding profits again to compensating variation generally double-counts them. A monetary equivalent is defined by an explicit transfer and price convention, with monotonicity and range conditions. Average effects, mechanism contributions, transfers and welfare losses are different targets. Infinite sums require convergence and dynamic feasible plans require terminal or no-Ponzi conditions.

\textbf{Source versions.} A working-paper date, revision date and journal publication year are distinguished. A DOI for a journal version is not automatically the DOI of an earlier manuscript. Locators refer to the stated version and printed pages where specified. A metadata-only check does not verify a claim about a full-text conclusion. Every entry's source subsection narrows the provenance claim accordingly.
% END ECONOMICS STATEMENT
\end{document}
